\section{Finite-prefix verification details}

This chapter describes how the finite-prefix base is verified in
Lean.  The emphasis is on the structure of the verification, not on a
transcription of the file.

\subsection{The exact expansion}

The finite base is the conjunction
\begin{equation*}
  \Orbit_7(0)=7,\ \Orbit_7(1)=9,\ \dots,\ \Orbit_7(16)=34177.
\end{equation*}
In Lean this is the statement
\texttt{fullOrbitIter\_\allowbreak prefix\_\allowbreak expand}.  The
proof proceeds by unfolding the orbit definition and rewriting each
step with the exact division
\begin{equation*}
  5\,\Orbit_7(n)+1=2^{t_n}\,\Orbit_7(n+1).
\end{equation*}

\subsection{Per-step table}

\begin{center}
\scriptsize
\begin{tabular}{@{}rrrrr@{}}
\toprule
$n$ & $\Orbit_7(n)$ & $5\Orbit_7(n)+1$ & $t_n$ & $\Orbit_7(n+1)$ \\
\midrule
0 & 7 & 36 & 2 & 9 \\
1 & 9 & 46 & 1 & 23 \\
2 & 23 & 116 & 2 & 29 \\
3 & 29 & 146 & 1 & 73 \\
4 & 73 & 366 & 1 & 183 \\
5 & 183 & 916 & 2 & 229 \\
6 & 229 & 1146 & 1 & 573 \\
7 & 573 & 2866 & 1 & 1433 \\
8 & 1433 & 7166 & 1 & 3583 \\
9 & 3583 & 17916 & 2 & 4479 \\
10 & 4479 & 22396 & 2 & 5599 \\
11 & 5599 & 27996 & 2 & 6999 \\
12 & 6999 & 34996 & 2 & 8749 \\
13 & 8749 & 43746 & 1 & 21873 \\
14 & 21873 & 109366 & 1 & 54683 \\
15 & 54683 & 273416 & 3 & 34177 \\
\bottomrule
\end{tabular}
\end{center}

Each row is an exact identity of the form
\begin{equation*}
  5x+1=2^t y.
\end{equation*}
The verification expands the left-hand side, evaluates the power of
two, and checks the quotient; no numerical search is involved.

\subsection{Prefix weights}

The exact prefix weights are
\begin{equation*}
  0,2,3,5,6,7,9,10,11,12,14,16,18,20,21,22,25.
\end{equation*}
They are the partial sums of the exact step weights
\begin{equation*}
  2,1,2,1,1,2,1,1,1,2,2,2,2,1,1,3.
\end{equation*}
The Lean statement
\texttt{fullOrbit\_\allowbreak prefix\_\allowbreak step\_\allowbreak
weights} records the weight word.

\subsection{First big step}

\begin{lemma}
For every $n<15$, $t_n\le2$, and $t_{15}=3$.
\end{lemma}

\begin{proof}
The table gives $t_n\le2$ for $n<15$ and $t_{15}=3$ directly.
\end{proof}

The Lean statement
\texttt{fullOrbit\_\allowbreak first\_\allowbreak t\_\allowbreak ge3\_
\allowbreak is\_\allowbreak exactly\_\allowbreak 3} packages both
claims:
\begin{verbatim}
theorem fullOrbit_first_t_ge3_is_exactly_3 :
    (forall n : Nat, n < 15 ->
      twoValuation (5 * fullOrbitIter n + 1) <= 2) /\
    twoValuation (5 * fullOrbitIter 15 + 1) = 3
\end{verbatim}

\subsection{Uniqueness}

The uniqueness statement is
\begin{verbatim}
theorem first_big_step_unique (n : Nat)
    (hsmall : forall m : Nat, m < n ->
      twoValuation (5 * fullOrbitIter m + 1) <= 2)
    (hbig : 3 <= twoValuation (5 * fullOrbitIter n + 1)) :
    n = 15 /\ twoValuation (5 * fullOrbitIter n + 1) = 3
\end{verbatim}
The proof is a case split on $n$: for $n<15$ the hypothesis
\texttt{hbig} contradicts the table, for $n=15$ the table gives the
weight, and for $n>15$ the depth-$15$ step would already be first.

\subsection{Why explicit rewriting}

The repository records the finite base with explicit rewriting rather
than \texttt{native\_decide}.  The purpose is to keep the verification
inside the kernel's reduction and rewrite machinery, so no
code-generation step is trusted.  The paper's audit searches the
main-chain files for \texttt{native\_decide}; the current count is
zero.

\subsection{Use in the exclusions}

The finite base is used in exactly two ways:
\begin{itemize}[leftmargin=*]
\item the $d=3$ family needs a first big step of weight $6$, which is
  excluded because the true first big step has weight $3$;
\item the $d\ge4$ branches need a first big step of weight at least
  $5$, also excluded by the same fact.
\end{itemize}
No later orbit value is used.

\subsection{Status}

\begin{center}
\small
\begin{tabular}{@{}p{0.56\textwidth}p{0.36\textwidth}@{}}
\toprule
Statement & Status \\
\midrule
\texttt{fullOrbitIter\_\allowbreak prefix\_\allowbreak expand} &
  compiled \\
\texttt{fullOrbit\_\allowbreak prefix\_\allowbreak step\_\allowbreak
weights} & compiled \\
\texttt{fullOrbit\_\allowbreak first\_\allowbreak t\_\allowbreak ge3\_
\allowbreak is\_\allowbreak exactly\_\allowbreak 3} & compiled \\
\texttt{first\_\allowbreak big\_\allowbreak step\_\allowbreak unique} &
  compiled \\
no \texttt{native\_decide} in the finite base & verified by inspection \\
\bottomrule
\end{tabular}
\end{center}
